\section{Methods and implementation}
\subsection{Interfaces and persistent execution}
Figure~\ref{fig:workflow} summarizes the separation of responsibilities.
The MD interface supplies positions, forces, species, stable atom identities,
cell information, and step indices. Segmented execution provides explicit
restart boundaries at which an accepted model can replace the active model.
Simulation state and model identity are recorded together so that a resumed
segment can be traced to the potential that generated it.

The controller persists its stage and records intended external actions,
their execution, and their outcomes. On restart, unresolved actions are
reconciled against the relevant backend before a new job is submitted.
Event, label, dataset, and model identities link the corresponding records.
The SQLite store uses write-ahead logging and uniqueness constraints.
These mechanisms support the recovery experiments reported below; their
presence alone is not evidence of correctness under every scheduler or
distributed failure mode.

Reference calculations are collected and validated before they enter a
dataset. Model promotion is separate from model training: a candidate must
pass the configured validation, export-consistency, and engine checks.
A rejected candidate leaves the accepted model available. The pilot uses
permissive promotion thresholds, stated in Section~\ref{sec:protocol};
passing those checks is not interpreted as proof of an accurate force field.

\subsection{Monitoring, events, and hotspot selection}
The monitoring interface supports force magnitude, force change,
nearest-neighbor distance, displacement, coordination change, local
Lennard--Jones (LJ) compatibility, and committee force dispersion.
Metrics are scheduled in a light stage, a physics stage, and a full stage.
For atom $i$ and stage $s$, the configured aggregate has the form
\begin{equation}
 S_i^{(s)}=\sum_{k\in\mathcal{K}_s}w_k g_k(x_{ik}),
 \label{eq:score}
\end{equation}
where $g_k$ is a metric-specific transformation and $w_k$ is a configured
weight. Several transformations use gated absolute standardized deviations
from running statistics. Other metrics retain raw-value contributions.
Unavailable metrics contribute zero and are excluded from the associated
statistics. Consequently, these aggregates are heuristic trigger scores,
not calibrated probabilities of reference-force error.

The pilot stage thresholds are 4, 5, and 6, with one atom sufficient to
trigger a stage. Under triggered scheduling, a lower-stage decision
controls whether the next stage is evaluated. The implementation can instead
evaluate stages unconditionally; the exploratory static-pool experiment
uses this option. The executed online pilot disables force-change,
displacement, coordination-change, and separate reference-model-deviation
metrics. It therefore does not validate every temporal feature available
in the platform.

For a committee of $M$ models, the force-dispersion feature is
\begin{equation}
 u_i=\left[\frac{1}{M}\sum_{m=1}^{M}
 \left\|\mathbf{F}^{(m)}_i-\overline{\mathbf{F}}_i\right\|^2\right]^{1/2},
 \qquad
 \overline{\mathbf{F}}_i=\frac{1}{M}\sum_{m=1}^{M}\mathbf{F}^{(m)}_i .
 \label{eq:committee}
\end{equation}
It is a population dispersion of force vectors, without an additional
factor of three.

The LJ feature fits the current predicted force to two local geometric
basis vectors. With $\mathbf{d}_{ij}=\mathbf{r}_j-\mathbf{r}_i$ and
$r_{ij}=\|\mathbf{d}_{ij}\|$,
\begin{align}
 \mathbf{a}_i&=-2\sum_j\frac{\mathbf{d}_{ij}}{r_{ij}^{14}},
 & \mathbf{b}_i&=\sum_j\frac{\mathbf{d}_{ij}}{r_{ij}^{8}},\\
 \widehat{\mathbf{F}}^{\rm LJ}_i&=\alpha_i\mathbf{a}_i+\beta_i\mathbf{b}_i,
 & \rho_i&=\frac{\|\widehat{\mathbf{F}}^{\rm LJ}_i-\mathbf{F}_i\|}
 {\|\mathbf{F}_i\|+\delta}.
 \label{eq:lj}
\end{align}
Here $\alpha=24\epsilon\sigma^{12}$ and
$\beta=24\epsilon\sigma^6$ for the conventional LJ potential, and the force
sign follows $\mathbf{F}=-\nabla E$. The pilot fits each atom separately
using neighbors within 6~\AA, including covalently bonded neighbors.
Rank, conditioning, positivity, and length-scale checks reject unsuitable
fits (Appendix~\ref{app:settings}). This is a compatibility feature for a
simple pair-force form. It is not a physicality certificate for a covalent
force field. Two fitted coefficients span a plane in three-dimensional
force space, and a positive uniform force rescaling can be absorbed into
the fitted energy scale.

Triggered frames are retained with configured pre- and post-trigger
context. Additional triggers within a pending event merge reasons,
scores, and hotspot identities; incomplete events are marked explicitly.
Trigger-frame identities are retained separately from later event members.
Hotspot atoms are clustered by a 3~\AA\ single-linkage distance.
Candidates are ranked by score, then filtered against selected and
previously labeled geometric fingerprints. The pilot uses a fingerprint
distance threshold of 0.1. These filters can determine the acquired
labels even when many frames trigger.

\subsection{Extraction, capping, and a fragment-fidelity gate}
The implementation offers graph, radial-cluster, block, and slab extraction.
The molecular experiments here assess graph and radial extraction only.
Whole carbon--hydrogen units are retained so that terminal hydrogens
follow their bonded carbon. Every severed carbon--carbon bond is
terminated by a hydrogen at 1.09~\AA\ along the cut-bond direction.
Caps are not relaxed in these experiments.

The calibration compares graph neighborhoods of two, three, and four
carbon--carbon bond hops with radial neighborhoods having a
1.5~\AA\ core and buffers of 3, 4, or 5~\AA.
For a perturbation site consisting of several atoms, the radial
neighborhood is the union of their neighborhoods; it is not necessarily
a sphere about a single atom. Parent identities, retained-atom identities,
region roles, and cap records are preserved.

Each candidate extraction setting is checked against parent DFT forces
computed with the same electronic-structure settings. Over the evaluated
core atoms $\mathcal{C}$, the gate measures the component root-mean-square
error (RMSE) and maximum atomwise vector error:
\begin{align}
 R_{\rm core}&=
 \left[\frac{1}{3|\mathcal{C}|}
 \sum_{i\in\mathcal{C}}\|\mathbf{F}^{\rm frag}_i-
                         \mathbf{F}^{\rm parent}_i\|^2\right]^{1/2},\\
 E_{\max}&=\max_{i\in\mathcal{C}}
 \|\mathbf{F}^{\rm frag}_i-\mathbf{F}^{\rm parent}_i\|.
 \label{eq:fidelity}
\end{align}
Acceptance requires $R_{\rm core}\leq0.03$~\forceunit,
$E_{\max}\leq0.10$~\forceunit, complete evaluation, and no failed
fragment. The lowest measured mean-cost passing setting is selected;
a fixed order resolves a cost tie. The settings and acceptance criteria
were recorded before the fidelity sweep. This local protocol record
is not an external preregistration.

\subsection{Force supervision and model updates}
Local datasets distinguish core, inner-buffer, outer-buffer, boundary,
and cap atoms. In the executed study, force weights are 1 for the core,
0.3 for the inner buffer, and zero for the outer buffer, boundary,
caps, and unassigned atoms. Full-system force labels have unit weights.
The weighted force root-mean-square loss is
\begin{equation}
 L_F=\left[
 \frac{\sum_i w_i\sum_{\mu\in\{x,y,z\}}
 (F_{i\mu}^{\theta}-F_{i\mu}^{\rm ref})^2}
 {3\sum_i w_i}\right]^{1/2}.
 \label{eq:loss}
\end{equation}
The implementation handles a zero-weight batch explicitly and uses a
numerical floor for a finite derivative at zero error. Zero-weight targets
contribute neither a force-loss term nor a force gradient.
Capped-fragment energies are excluded from energy supervision: a fragment
energy is not a parent energy or an additive parent-energy contribution.
Full-system labels retain their own energy targets.

Local DFT fidelity and learned-model context are separate constraints.
Equation~\eqref{eq:fidelity} tests whether the reference forces survive
fragmentation. It does not test whether the trained model predicts the
same core force on the fragment and its parent. We report the completed
model-context diagnostic as a limitation rather than treating a DFT gate
pass as evidence that the learned representation is insensitive to
removed neighbors.
